\newif\ifglobal

%%%%%%%%%%%%%%%%%%%%%%%
%%% Compile to view %%%
%%%%%%%%%%%%%%%%%%%%%%%

%%%%%%%%%%%%%%%%%%%%%%%%%%%%%%%%%%%%%%%%%%%%%%%%%%%
%%% Readme:                                     %%%
%%% https://www.overleaf.com/read/bwdjvfjksvjy/ %%%
%%%%%%%%%%%%%%%%%%%%%%%%%%%%%%%%%%%%%%%%%%%%%%%%%%%

% >>> M?: marks a module setting number or important notice

% >>> M4: Set {./relative-path} to external files for this module <<<
\def \modulefiles {./30-LEDPWM}
% To add an external file, use:
% (Replace '---' with actual filename)
% '\input{\modulefiles/tables/---.tex}' for tables
% '\includegraphics{\modulefiles/figures/---}' for figures
% '\subfile{\modulefiles/---}' for register files

% >>> M1: This module is in global project? <<<
% 'YES' - uncomment; 'NO' - comment out
\globaltrue

\ifglobal
    \documentclass[main\_\_CN.tex]{subfiles}

\else
    % Font "FandolSong-Regular" used by ctex does not have GSUB table.
    % It causes warnings that do not affect compilation and can be ignored.
    \PassOptionsToPackage{quiet}{fontspec}

    \documentclass[openany, 10pt]{book}

    % >>> M2: In ./chip-spec-settings/preamble-trm-module, also find
    %         settings for visibility of todo notes, watermark <<<
    \usepackage{preamble-trm-module}


    % Enable Chinese typesetting
    \usepackage{ctex}

    % >>> M3: Temporary labels in Chinese
    %         you can add your own labels there <<<
    \myexternaldocument{temp-labels-cn}

\fi %\ifglobal


\setcounter{secnumdepth}{4}

% >>> M5: If this module needs any updates in
% './chip-spec-settings/preamble-trm-module.sty'
% (include additional package, etc.), contact your module PM
% or create the file './chip-spec-settings/changelog.tex' make the
% necessary updates yourself and document them in this file <<<

\raggedright
\raggedbottom


\begin{document}


% Sets language into which pre-defined bits of text are translated
\selectlanguage{chinese}
\setchineseformatting

% Do not delete this block
\ifglobal
% Add the following front matter if
% this module is outside of global project
\else
    \InputIfFileExists{readme}{}{}
    {\let\clearpage\relax \listoftodos}
    {\let\clearpage\relax \listofdonetodos}
    \tableofcontents
    %\listoftables
    %\listoffigures
\fi


%%%%%%%%%%%%%%%%%%%%%%%%%
%%% TRM Module Begins %%%
%%%%%%%%%%%%%%%%%%%%%%%%%

\hypertarget{ledpwm}{}
\chapter{LED PWM 控制器 (LEDC)}
\label{mod:ledpwm}

\section{概述}

LED PWM 控制器主要用于控制 LED 设备，也可以产生 PWM 信号用于控制其他开关设备。该控制器由 14 位定时器和波形发生器组成。

\section{特性}
LED PWM 控制器具有如下特性:
\begin{itemize}
    \item 四个独立定时器，可实现小数分频
    \item 八个独立波形发生器，可产生 8 路 PWM 信号
    \item PWM 信号占空比可递增或递减渐变，无须处理器干预。渐变完成可产生中断事件
    \item 输出 PWM 信号相位可调
    \item 低功耗模式下可输出 PWM 信号

\end{itemize}

为方便描述，下文中八个 PWM 发生器统称为 PWM\regindex{n}，四个定时器统称为 Timer\regindex{x}。

\section{功能描述}

\subsection{架构}

图 \ref{Figure:led_pwm_architecture}为 LED PWM 控制器的架构。图 \ref{Figure:led_pwm_PWM generator} 为一个 PWM 生成器及其选取的定时器和计数器。

\begin{figure}[ht]
    \centering
    \includegraphics[width=0.6\textwidth]{\modulefiles/figures/led2-1_723}
    \caption{LED\_PWM 架构}
    \label{Figure:led_pwm_architecture}
\end{figure}
\begin{figure}[ht]
    \centering
    \includegraphics[width=1.0\textwidth]{\modulefiles/figures/ledc2-2}
    \caption{LED\_PWM 生成器图}
    \label{Figure:led_pwm_PWM generator}
\end{figure}
\subsection{定时器}\label{ledc_timers}

\begin{figure}[ht]
    \centering
    \fbox{\includegraphics[width=0.9\textwidth]{\modulefiles/figures/ledc_frag_div}}
    \caption{LED\_PWM 分频器}
    \label{Figure:led_pwm_divider}
\end{figure}

LED PWM 控制器的时钟 LEDC\_PWM\_CLK 有三个时钟源：APB\_CLK、RC\_FAST\_CLK 和 XTAL\_CLK，可通过配置 \hyperref[fielddesc:LEDCAPBCLKSEL]{LEDC\_APB\_CLK\_SEL}[1:0] 来选择。LED PWM 定时器的时钟 LEDC\_CLK\regindex{x} 有两个时钟源：LEDC\_PWM\_CLK 和 REF\_TICK。如使用 REF\_TICK 作定时器的时钟源，需将 \hyperref[fielddesc:LEDCAPBCLKSEL]{LEDC\_APB\_CLK\_SEL}[1:0] 置 1，REF\_TICK 的周期需为 APB\_CLK 周期的整数倍，否则定时器的时钟会不精准。更多关于时钟源的信息请参考 \hyperref[mod:resclk]{复位与时钟}一章。

LEDC\_CLK\regindex{x} 经分频器分频后的输出时钟用作计数器的基准时钟。分频器的分频系数通过 \hyperref[fielddesc:LEDCCLKDIVTIMERX]{LEDC\_CLK\_DIV\_TIMER\regindex{x}} 寄存器配置。该分频系数为定点数，高 10 位为整数部分 A，低 8 位为小数部分 B。分频系数 LEDC\_CLK\_DIV\regindex{x} 的公式为：

\begin{center}
$LEDC\_CLK\_DIV\regindex{x} = A + \frac{B}{256}$
\end{center}

小数部分 B 不为 0 时，分频器的输入和输出时钟如图 \ref{Figure:led_pwm_divider}所示。256 个输出时钟中有 B 个以 (A+1) 分频，有 (256-B) 个以 A 分频。以 (A+1) 分频的 B 个输出时钟均匀分布在 256 个输出时钟中。

LED PWM 控制器的计数器为 14 位，计数最大值为 $2^{\hyperref[fielddesc:LEDCTIMERXDUTYRES]{LEDC\_TIMER\regindex{x}\_DUTY\_RES}}-1$。每当计数值达到\\ $2^{\hyperref[fielddesc:LEDCTIMERXDUTYRES]{LEDC\_TIMER\regindex{x}\_DUTY\_RES}}-1$ 时，计数器便会溢出，并重新从 0 开始计数。软件可以复位、暂停并读取计数器的计数值。计数器每次溢出时可产生 LEDC\_TIMER\regindex{x}\_OVF\_INT 中断，也可在溢出 (\hyperref[fielddesc:LEDCOVFNUMCHN]{LEDC\_OVF\_NUM\_CH\regindex{n}} + 1) 次时产生中断。寄存器配置流程如下：


 \begin{enumerate}
    \item 置位 \hyperref[fielddesc:LEDCOVFCNTENCHN]{LEDC\_OVF\_CNT\_EN\_CH\regindex{n}}
    \item 将 \hyperref[fielddesc:LEDCOVFNUMCHN]{LEDC\_OVF\_NUM\_CH\regindex{n}} 配置为溢出次数减 1
    \item 置位\hyperref[fielddesc:LEDCOVFCNTCHNINTENA]{LEDC\_OVF\_CNT\_CH\regindex{n}\_INT\_ENA}
    \item 置位 \hyperref[fielddesc:LEDCTIMERXDUTYRES]{LEDC\_TIMER\regindex{x}\_DUTY\_RES} 启动定时器，等待 LEDC\_OVF\_CNT\_CH\regindex{n}\_INT 中断
\end{enumerate}

PWM 生成器输出信号 sig\_out\regindex{n} 的频率取决于定时器时钟源 LEDC\_CLK\regindex{x} 的频率、 时钟分频系数 LEDC\_CLK\_DIV\regindex{x} 以及占空比精度（计数器位宽）hyperref[fielddesc:LEDCTIMERXDUTYRES]{LEDC\_TIMER\regindex{x}\_DUTY\_RES}：

PWM 生成器输出信号 sig\_out\regindex{n} 的频率取决于分频器的分频系数以及计数器的计数范围：
    \[
    f_{\textrm{sig\_out\regindex{n}}}=\frac{f_{\textrm{LEDC\_CLK\regindex{x}}}}
                        {\textrm{LEDC\_CLK\_DIV}\regindex{x}\cdot 2^{hyperref[fielddesc:LEDCTIMERXDUTYRES]{LEDC\_TIMER\regindex{x}\_DUTY\_RES}}}
    \]

上述公式变形后，可得到以下公式计算预期的占空比精度：

    \[
    \textrm{hyperref[fielddesc:LEDCTIMERXDUTYRES]{LEDC\_TIMER\regindex{x}\_DUTY\_RES}}
    =\log_2\left(\frac{f_{\textrm{LEDC\_CLK\regindex{x}}}}
    {f_{\textrm{sig\_out\regindex{n}}}\cdot \textrm{LEDC\_CLK\_DIV\regindex{x}}}\right)
    \]

表 \ref{tab:ledc-common-freq-res} 列出了常用配置频率及其对应精度。

\begin{ThreePartTable}

    \begin{TableNotes}

        \item[1] 最高精度指时钟分频系数 LEDC\_CLK\_DIV\regindex{x} 为 1 时的精度，向下取整。如果经公式计算出的最高精度超过了计数器位宽 14 位，则最高精度为 14。

        \item[2] 最低精度指时钟分频系数 LEDC\_CLK\_DIV\regindex{x} 为 $1023+\frac{255}{256}$ 时的精度，向上取整。如果经公式计算出的最低精度小于 0，则最低精度为 1。

    \end{TableNotes}

\begin{longtable}[c]{ | l | l | l | l | }
\caption{常用配置频率及精度}
\label{tab:ledc-common-freq-res}
\\\hline
\rowcolor{lightgray}
LEDC\_CLK\regindex{x} & PWM 频率 & 最高精度（位） \tnote{1} & 最低精度（位）\tnote{2} \\\hline

\endhead
\insertTableNotes
\endfoot

APB\_CLK (80 MHz) & 1 kHz & 14 & 7 \\\hline
APB\_CLK (80 MHz) & 5 kHz & 13 & 4 \\\hline
APB\_CLK (80 MHz) & 10 kHz & 12 & 3 \\\hline
RC\_FAST\_CLK (8 MHz) & 1 kHz & 12 & 3 \\\hline
RC\_FAST\_CLK (8 MHz) & 2 kHz & 11 & 2 \\\hline
REF\_TICK (1 MHz) & 1 kHz & 9 & 1 \\\hline
\end{longtable}
\end{ThreePartTable}


重新设置分频器的分频系数和计数器的溢出值，需分别配置 \hyperref[fielddesc:LEDCCLKDIVTIMERX]{LEDC\_CLK\_DIV\_TIMER\regindex{x}} 和 \hyperref[fielddesc:LEDCTIMERXDUTYRES]{LEDC\_TIMER\regindex{x}\_DUTY\_RES} ，然后置位 \hyperref[fielddesc:LEDCTIMERXPARAUP]{LEDC\_TIMER\regindex{x}\_PARA\_UP}，否则配置无效。新的配置在计数器溢出时更新。

\subsection{PWM 生成器}\label{LEDCPWM generator}

如图 \ref{Figure:led_pwm_PWM generator}所示，每个 PWM 生成器主要包括一个高低电平比较器和两个选择器。PWM 生成器获取 LED PWM 定时器的 14 位计数值，与高低电平比较器的值 Hpoint\regindex{n} 和 Lpoint\regindex{n} 比较，来控制 PWM 信号的电平。

\begin{itemize}
\item 如果 Timer\regindex{x}\_cnt == Hpoint\regindex{n}，则 sig\_out\regindex{n} 为 1。
\item 如果 Timer\regindex{x}\_cnt == Lpoint\regindex{n}，则 sig\_out\regindex{n} 为 0。
\end{itemize}

Hpoint\regindex{n} 由 \hyperref[fielddesc:LEDCHPOINTCHN]{LEDC\_HPOINT\_CH\regindex{n}} 在计数器溢出时更新。Lpoint\regindex{n} 的初始值为计数器溢出时 \hyperref[fielddesc:LEDCDUTYCHN]{LEDC\_DUTY\_CH\regindex{n}}[18..4] 和 \hyperref[fielddesc:LEDCHPOINTCHN]{LEDC\_HPOINT\_CH\regindex{n}} 的和。通过配置以上两个字段，可设置 PWM 输出的相对相位和占空比。

图 \ref{Figure:led_pwm_fixed_duty} 为占空比固定时的 PWM 波形。

\begin{figure}[ht]
    \centering
    \fbox{\includegraphics[width=0.6\textwidth]{\modulefiles/figures/ledc_2-3-1_723}}
    \caption{LED\_PWM 输出信号图}
    \label{Figure:led_pwm_fixed_duty}
\end{figure}

\label{ledcfd:dithering}
\hyperref[fielddesc:LEDCDUTYCHN]{LEDC\_DUTY\_CH\regindex{n}} 是一个具有四位小数的定点寄存器。\hyperref[fielddesc:LEDCDUTYCHN]{LEDC\_DUTY\_CH\regindex{n}}[18..4] 为整数部分，可直接用于 PWM 计算。\hyperref[fielddesc:LEDCDUTYCHN]{LEDC\_DUTY\_CH\regindex{n}}[3..0] 为小数部分，可用于调整输出信号。如 \hyperref[fielddesc:LEDCDUTYCHN]{LEDC\_DUTY\_CH\regindex{n}}[3..0] 不为 0，那么 sig\_out\regindex{n} 每 16 个周期中，有 \hyperref[fielddesc:LEDCDUTYCHN]{LEDC\_DUTY\_CH\regindex{n}}[3..0] 个周期的 PWM 脉冲宽度要比 (16 - \hyperref[fielddesc:LEDCDUTYCHN]{LEDC\_DUTY\_CH\regindex{n}}[3..0]) 个周期的脉冲宽度多一个定时器的计数周期。该功能可将 PWM 生成器的精度提至 18 位。

\subsection{渐变占空比}

LED PWM 输出信号可由一种占空比渐变为另一种占空比。渐变功能由
\hyperref[fielddesc:LEDCDUTYCHN]{LEDC\_DUTY\_CH\regindex{n}}、\hyperref[fielddesc:LEDCDUTYSTARTCHN]{LEDC\_DUTY\_START\_CH\regindex{n}}、\hyperref[fielddesc:LEDCDUTYINCCHN]{LEDC\_DUTY\_INC\_CH\regindex{n}} 、\hyperref[fielddesc:LEDCDUTYNUMCHN]{LEDC\_DUTY\_NUM\_CH\regindex{n}} 和 \hyperref[fielddesc:LEDCDUTYSCALECHN]{LEDC\_DUTY\_SCALE\_CH\regindex{n}} 配置。

\hyperref[fielddesc:LEDCDUTYSTARTCHN]{LEDC\_DUTY\_START\_CH\regindex{n}} 用于更新 Lpoint\regindex{n} 的值。如置位该位，计数器溢出时 Lpoint\regindex{n} 会自动递增或递减。\\\hyperref[fielddesc:LEDCDUTYINCCHN]{LEDC\_DUTY\_INC\_CH\regindex{n}} 位用于设置渐变模式。

sig\_out\regindex{n} 的占空比每隔 \hyperref[fielddesc:LEDCDUTYCYCLECHN]{LEDC\_DUTY\_CYCLE\_CH\regindex{n}} 个 PWM 脉冲周期，便会在当前占空比上加上或减去\\ \hyperref[fielddesc:LEDCDUTYSCALECHN]{LEDC\_DUTY\_SCALE\_CH\regindex{n}}。

占空比渐变次数达到 \hyperref[fielddesc:LEDCDUTYNUMCHN]{LEDC\_DUTY\_NUM\_CH\regindex{n}} 时停止渐变，并产生 LEDC\_DUTY\_CHNG\_END\_CH\regindex{n}\_INT。图 \ref{Figure:led_pwm_unfixed_duty}为渐变的图例。该配置下，每 \hyperref[fielddesc:LEDCDUTYCYCLECHN]{LEDC\_DUTY\_CYCLE\_CH\regindex{n}} 个 PWM 脉冲周期，sig\_out\regindex{n}的占空比会增加 \hyperref[fielddesc:LEDCDUTYSCALECHN]{LEDC\_DUTY\_SCALE\_CH\regindex{n}}。

\begin{figure}[ht]
    \centering
    \fbox{\includegraphics[width=0.8\textwidth]{\modulefiles/figures/ledc_f5}}
    \caption{渐变占空比输出信号图}
    \label{Figure:led_pwm_unfixed_duty}
\end{figure}

\hyperref[fielddesc:LEDCSIGOUTENCHN]{LEDC\_SIG\_OUT\_EN\_CH\regindex{n}} 用于启动 PWM 波形输出。\hyperref[fielddesc:LEDCSIGOUTENCHN]{LEDC\_SIG\_OUT\_EN\_CH\regindex{n}} 为 0 时，sig\_out\regindex{n}的电平值恒定为 \hyperref[fielddesc:LEDCIDLELVCHN]{LEDC\_IDLE\_LV\_CH\regindex{n}}。

软件更新 \hyperref[fielddesc:LEDCHPOINTCHN]{LEDC\_HPOINT\_CH\regindex{n}}、\hyperref[fielddesc:LEDCDUTYSTARTCHN]{LEDC\_DUTY\_START\_CH\regindex{n}}、\hyperref[fielddesc:LEDCSIGOUTENCHN]{LEDC\_SIG\_OUT\_EN\_CH\regindex{n}}、 \hyperref[fielddesc:LEDCTIMERSELCHN]{LEDC\_TIMER\_SEL\_CH\regindex{n}}、\\\hyperref[fielddesc:LEDCDUTYNUMCHN]{LEDC\_DUTY\_NUM\_CH\regindex{n}}、 \hyperref[fielddesc:LEDCDUTYCYCLECHN]{LEDC\_DUTY\_CYCLE\_CH\regindex{n}}、 \hyperref[fielddesc:LEDCDUTYSCALECHN]{LEDC\_DUTY\_SCALE\_CH\regindex{n}}、\hyperref[fielddesc:LEDCDUTYINCCHN]{LEDC\_DUTY\_INC\_CH\regindex{n}} \\和 \hyperref[fielddesc:LEDCOVFCNTENCHN]{LEDC\_OVF\_CNT\_EN\_CH\regindex{n}} 的配置后，需置位 \hyperref[fielddesc:LEDCPARAUPCHN]{LEDC\_PARA\_UP\_CH\regindex{n}} 应用新配置。

\subsection{中断}
\begin{itemize}
\item \label{int:LEDCOVFCNTCHnINT}LEDC\_OVF\_CNT\_CH\regindex{n}\_INT：定时器计数器溢出 (\hyperref[fielddesc:LEDCOVFNUMCHN]{LEDC\_OVF\_NUM\_CH\regindex{n}} + 1) 次且寄存器 \hyperref[fielddesc:LEDCOVFCNTENCHN]{LEDC\_OVF\_CNT\_EN\_CH\regindex{n}} 置 1 时触发中断。
\item \label{int:LEDCDUTYCHNGENDCHnINT}LEDC\_DUTY\_CHNG\_END\_CH\regindex{n}\_INT：LED PWM 生成器渐变完成后触发中断。
\item \label{int:LEDCTIMERxOVFINT}LEDC\_TIMER\regindex{x}\_OVF\_INT：LED PWM 定时器达到最大计数值时触发中断。
\end{itemize}

\section{基地址}\label{subsection:ledc base address}
用户可以通过两个不同的寄存器基地址访问 LED PWM，如表 \ref{table:ledc base address} 所示。更多信息，请访问章节 \ref{mod:sysmem} \textit{\nameref{mod:sysmem}}。

\begin{longtable}[c]{ | l |c| }
    \caption{LED PWM 控制器基地址} \label{table:ledc base address}
    \\\hline
    \rowcolor{lightgray}
    访问总线  &  基地址\\
    \hline
    PeriBUS1 & \hexprefix{}3F419000 \\\hline
    PeriBUS2 & \hexprefix{}60019000 \\\hline
\end{longtable}

\hypertarget{ledpwm-reg-summ}{}
\section{寄存器列表}

请注意，下表的地址指相对于 LED PWM 基地址的偏移量（相对地址）。请参阅章节 \ref{subsection:ledc base address} 获取有关 LED PWM 基地址的信息。请注意，下表中的地址都是相对于 LED PWM 基地址的地址偏移量（相对地址）。更多有关 LED PWM 基地址的信息，请前往 \ref{subsection:ledc base address} 章节。
\subfile{\modulefiles/ledc_summary_cn.tex}

\section{寄存器}
\subfile{\modulefiles/ledc_reg_cn.tex}

\end{document}
